\section{Independent review of Drinfeld and Hilbert-scheme applications}
\label{r20audit:DUHMroot}

\paragraph{Exact versions and verdict.}
The reviewed sources are
\path{families/drinfeld_uniformization_with_good_coefficients.tex},
SHA--256
\texttt{aa620c1d569aa9b52555fd492ea2511e0f544ec07fc52df6df9bd6fa3936068a},
and
\path{families/hilbert_schemes_and_actual_chow_transfer.tex},
SHA--256
\texttt{5170dcd3804d1adb44d5c72300cefbfb6756ac4785bb852ff68eb60fc6fa4f39}.
Both receive \textbf{PASS relative to their explicitly stated inputs}.
DP20's independent review is now complete. These are internal proof
reviews, not external peer review or priority certification.

\subsection{DU20}
I checked Ito's primary paper, Theorem 1.2 and Sections 6.1--6.2:
the projective algebraization and the finite-index strictly semistable
cover are genuine geometric inputs, not consequences inferred solely
from weight--monodromy. The closed strata are products of the
successive linear-space blowups $B_d$.

Each $B_d$ has a native split Tate Chow motive: its centers are the
strict transforms of rational linear spaces, lower-dimensional $B_b$,
with earlier intersections already blown up. The actual smooth blowup
maps give the motive decomposition after every test product.
Products and projective bundles preserve this property. The
all-stratum staircase formulas of AT20 therefore preserve it for any
finite product of the strict models, in arbitrary dimensions.

For a split Tate Chow motive, the complementary intersection pairing
on Chow groups is perfect; hence numerical zero is already zero in
Chow. This applies to every entry of PI20's error. The same rational
inverse consequently works at every invertible prime. Tensoring with
the arbitrary native good coefficient preserves the full operator
identity and does not require comparison of equivalences on that
coefficient.

The inertia argument uses Ito's all-degree theorem, purity for the
good coefficient, and the tensor rule for centered monodromy
filtrations. The tensor rule follows from the decomposition of
nilpotent strings, with weights added on graded tensor factors.
All dimensions at most two are covered separately by frozen endpoints,
so AK19 is not invoked beyond its stated range.

Every unramified extension retains the strict models and actual Tate
decompositions and permits new cycles. RI18 then supplies the asserted
all-finite extension statement. Finite etale descent of the cover uses
its complete degree; the complement is exactly the inverse image of
the original support. No nonfinite alteration pushforward of actual
forms occurs. The arbitrary coefficient is required to be good over
the native supplied model field. The source correctly does not infer
this for an arbitrary potentially-good coefficient from WMC alone.
The global statement is a placewise local criterion, with a common
field at each place and full-degree compositum descent.

\subsection{HM20}
I checked de Cataldo--Migliorini's primary author PDF, Theorem 6.2.1,
Proposition 6.1.5, and Remark 6.2.2. The last explicitly allows any
ground field; the relevant identity is an actual rational Chow-motive
identity, not merely a complex cohomological decomposition.
The finite group average on $S^{l(\lambda)}$ is an actual idempotent,
so singular symmetric products are represented by projectors on smooth
powers and are never treated as smooth analytic Green targets.

For a partition of $m$ of length $l$, the incidence has dimension
$m+l$ inside $S^l\times S^{[m]}$. Its transpose sends codimension $r$
to $r-(m-l)$, while the forward correspondence adds $m-l$. The cited
normalization constants are part of the actual theorem, and the
composites sum to the actual Chow diagonal. Tensoring these identities
with all other diagonals gives actual finite factorization through
the smooth products covered by DP20.

The frozen CV16 covariance is used only on the obstruction kernel.
Geometric homological triviality is preserved by the actual
correspondences. Impossible degrees vanish; degree-zero homologically
trivial cycles are zero componentwise. Thus the diagonal identity
forces the original obstruction to vanish. VT17 constructs the final
form for the specified original cycle, independently of supports of
the intermediate correspondence images.

The correspondence identities base change, and DP20 is reapplied to
new cycles after each finite extension. The height statement includes
both full complementary Chow groups and requires potential good
reduction of the extra factor at every finite place. No direct
semistable model of a Hilbert scheme, finite-dimensional Chow-motive
conjecture, positivity theorem, or arbitrary-variety result is used.
The R9 family candidates and prior Drinfeld literature are explicitly
acknowledged; no first-discovery claim is certified.

